% Options for packages loaded elsewhere
% Options for packages loaded elsewhere
\PassOptionsToPackage{unicode}{hyperref}
\PassOptionsToPackage{hyphens}{url}
\PassOptionsToPackage{dvipsnames,svgnames,x11names}{xcolor}
%
\documentclass[
  11pt,
  letterpaper,
  DIV=11,
  numbers=noendperiod]{scrartcl}
\usepackage{xcolor}
\usepackage[margin=2.5cm]{geometry}
\usepackage{amsmath,amssymb}
\setcounter{secnumdepth}{5}
\usepackage{iftex}
\ifPDFTeX
  \usepackage[T1]{fontenc}
  \usepackage[utf8]{inputenc}
  \usepackage{textcomp} % provide euro and other symbols
\else % if luatex or xetex
  \usepackage{unicode-math} % this also loads fontspec
  \defaultfontfeatures{Scale=MatchLowercase}
  \defaultfontfeatures[\rmfamily]{Ligatures=TeX,Scale=1}
\fi
\usepackage{lmodern}
\ifPDFTeX\else
  % xetex/luatex font selection
\fi
% Use upquote if available, for straight quotes in verbatim environments
\IfFileExists{upquote.sty}{\usepackage{upquote}}{}
\IfFileExists{microtype.sty}{% use microtype if available
  \usepackage[]{microtype}
  \UseMicrotypeSet[protrusion]{basicmath} % disable protrusion for tt fonts
}{}
\makeatletter
\@ifundefined{KOMAClassName}{% if non-KOMA class
  \IfFileExists{parskip.sty}{%
    \usepackage{parskip}
  }{% else
    \setlength{\parindent}{0pt}
    \setlength{\parskip}{6pt plus 2pt minus 1pt}}
}{% if KOMA class
  \KOMAoptions{parskip=half}}
\makeatother
% Make \paragraph and \subparagraph free-standing
\makeatletter
\ifx\paragraph\undefined\else
  \let\oldparagraph\paragraph
  \renewcommand{\paragraph}{
    \@ifstar
      \xxxParagraphStar
      \xxxParagraphNoStar
  }
  \newcommand{\xxxParagraphStar}[1]{\oldparagraph*{#1}\mbox{}}
  \newcommand{\xxxParagraphNoStar}[1]{\oldparagraph{#1}\mbox{}}
\fi
\ifx\subparagraph\undefined\else
  \let\oldsubparagraph\subparagraph
  \renewcommand{\subparagraph}{
    \@ifstar
      \xxxSubParagraphStar
      \xxxSubParagraphNoStar
  }
  \newcommand{\xxxSubParagraphStar}[1]{\oldsubparagraph*{#1}\mbox{}}
  \newcommand{\xxxSubParagraphNoStar}[1]{\oldsubparagraph{#1}\mbox{}}
\fi
\makeatother


\usepackage{longtable,booktabs,array}
\newcounter{none} % for unnumbered tables
\usepackage{calc} % for calculating minipage widths
% Correct order of tables after \paragraph or \subparagraph
\usepackage{etoolbox}
\makeatletter
\patchcmd\longtable{\par}{\if@noskipsec\mbox{}\fi\par}{}{}
\makeatother
% Allow footnotes in longtable head/foot
\IfFileExists{footnotehyper.sty}{\usepackage{footnotehyper}}{\usepackage{footnote}}
\makesavenoteenv{longtable}
\usepackage{graphicx}
\makeatletter
\newsavebox\pandoc@box
\newcommand*\pandocbounded[1]{% scales image to fit in text height/width
  \sbox\pandoc@box{#1}%
  \Gscale@div\@tempa{\textheight}{\dimexpr\ht\pandoc@box+\dp\pandoc@box\relax}%
  \Gscale@div\@tempb{\linewidth}{\wd\pandoc@box}%
  \ifdim\@tempb\p@<\@tempa\p@\let\@tempa\@tempb\fi% select the smaller of both
  \ifdim\@tempa\p@<\p@\scalebox{\@tempa}{\usebox\pandoc@box}%
  \else\usebox{\pandoc@box}%
  \fi%
}
% Set default figure placement to htbp
\def\fps@figure{htbp}
\makeatother





\setlength{\emergencystretch}{3em} % prevent overfull lines

\providecommand{\tightlist}{%
  \setlength{\itemsep}{0pt}\setlength{\parskip}{0pt}}



 


\KOMAoption{captions}{tableheading}
\makeatletter
\@ifpackageloaded{caption}{}{\usepackage{caption}}
\AtBeginDocument{%
\ifdefined\contentsname
  \renewcommand*\contentsname{Table of contents}
\else
  \newcommand\contentsname{Table of contents}
\fi
\ifdefined\listfigurename
  \renewcommand*\listfigurename{List of Figures}
\else
  \newcommand\listfigurename{List of Figures}
\fi
\ifdefined\listtablename
  \renewcommand*\listtablename{List of Tables}
\else
  \newcommand\listtablename{List of Tables}
\fi
\ifdefined\figurename
  \renewcommand*\figurename{Figure}
\else
  \newcommand\figurename{Figure}
\fi
\ifdefined\tablename
  \renewcommand*\tablename{Table}
\else
  \newcommand\tablename{Table}
\fi
}
\@ifpackageloaded{float}{}{\usepackage{float}}
\floatstyle{ruled}
\@ifundefined{c@chapter}{\newfloat{codelisting}{h}{lop}}{\newfloat{codelisting}{h}{lop}[chapter]}
\floatname{codelisting}{Listing}
\newcommand*\listoflistings{\listof{codelisting}{List of Listings}}
\makeatother
\makeatletter
\makeatother
\makeatletter
\@ifpackageloaded{caption}{}{\usepackage{caption}}
\@ifpackageloaded{subcaption}{}{\usepackage{subcaption}}
\makeatother
\usepackage{bookmark}
\IfFileExists{xurl.sty}{\usepackage{xurl}}{} % add URL line breaks if available
\urlstyle{same}
\hypersetup{
  pdftitle={Análisis de comportamiento de usuarios en ecommerce},
  pdfauthor={Camila Durand; Mateo Troncoso},
  colorlinks=true,
  linkcolor={blue},
  filecolor={Maroon},
  citecolor={Blue},
  urlcolor={Blue},
  pdfcreator={LaTeX via pandoc}}


\title{Análisis de comportamiento de usuarios en ecommerce}
\usepackage{etoolbox}
\makeatletter
\providecommand{\subtitle}[1]{% add subtitle to \maketitle
  \apptocmd{\@title}{\par {\large #1 \par}}{}{}
}
\makeatother
\subtitle{Trabajo Práctico Final - Minería de Datos}
\author{Camila Durand \and Mateo Troncoso}
\date{2026-06-08}
\begin{document}
\maketitle

\renewcommand*\contentsname{Table of contents}
{
\setcounter{tocdepth}{3}
\tableofcontents
}

\section{Introducción}\label{introducciuxf3n}

En la actualidad, las plataformas de comercio electrónico generan
grandes volúmenes de datos relacionados con la navegación e interacción
de los usuarios. El análisis de estos datos permite comprender patrones
de comportamiento, preferencias de consumo y recorridos de exploración
dentro de una tienda online.

La minería de datos aplicada al comercio electrónico permite identificar
relaciones entre productos, descubrir patrones frecuentes de navegación
y generar conocimiento útil para estrategias de marketing, sistemas de
recomendación y técnicas de cross-selling.

En el presente trabajo se analizará un conjunto de datos perteneciente a
una tienda online de ropa correspondiente al año 2008. A partir de
técnicas de análisis exploratorio, visualización de datos, minería de
itemsets frecuentes, reglas de asociación y análisis secuencial, se
buscará comprender cómo interactúan los usuarios con los productos y
categorías del sitio web.

Además, se estudiarán diferencias de comportamiento entre distintos
países y se analizará la influencia de variables como ubicación visual,
categoría y rango de precio en la exploración de productos.

\section{Objetivos}\label{objetivos}

\subsection{Objetivo General}\label{objetivo-general}

Descubrir patrones de comportamiento de navegación dentro de una tienda
online de ropa con el fin de generar conocimiento útil para sistemas de
recomendación, estrategias de cross-selling y análisis de comportamiento
de usuarios.

\subsection{Objetivos Específicos}\label{objetivos-especuxedficos}

\begin{enumerate}
\def\labelenumi{\arabic{enumi}.}
\item
  Identificar productos que suelen visualizarse conjuntamente
\item
  Detectar combinaciones de productos relevantes para campañas de
  cross-selling
\item
  Comparar hábitos de navegación entre países
\item
  Encontrar secuencias típicas de exploración de productos
\item
  Generar conocimiento para sistemas de recomendación
\end{enumerate}

\subsection{Origen y Características}\label{origen-y-caracteruxedsticas}

El dataset proviene del repositorio UCI y contiene información de
navegación de una tienda online polaca de ropa para embarazadas durante
5 meses del año 2008.

{\def\LTcaptype{none} % do not increment counter
\begin{longtable}[]{@{}ll@{}}
\toprule\noalign{}
Característica & Valor \\
\midrule\noalign{}
\endhead
\bottomrule\noalign{}
\endlastfoot
Período & Abril - Agosto 2008 \\
Registros totales & 165,474 \\
Sesiones únicas & 24,026 \\
Productos distintos & 217 \\
Países representados & 42 \\
\end{longtable}
}

\section{Exploración y Preparación de
Datos}\label{exploraciuxf3n-y-preparaciuxf3n-de-datos}

\subsection{Limpieza Inicial}\label{limpieza-inicial}

\textbf{\hfill\break
Decisiones de limpieza:}

\begin{itemize}
\item
  Se eliminó la variable \texttt{year} porque todos los registros son de
  2008.
\item
  Se filtraron los códigos de dominio de red 43 al 47 de la variable
  \texttt{country\ ya\ que\ no\ representan\ a\ ningun\ país.}
\item
  Se transformaron variables numéricas a factores a etiquetas
  descriptivas.
\item
  La etiqueta del id de session la pasamos de tipo numérico a factor
  para implementar de manera efectiva los algoritmos correspondientes.
\end{itemize}

\subsection{Estadísticas
Descriptivas}\label{estaduxedsticas-descriptivas}

{\def\LTcaptype{none} % do not increment counter
\begin{longtable}[]{@{}lll@{}}
\toprule\noalign{}
Estadístico & Clicks por sesión & Precio (USD) \\
\midrule\noalign{}
\endhead
\bottomrule\noalign{}
\endlastfoot
Mínimo & 1 & 18 \\
Mediana & 6 & 43 \\
Media & 9.8 & 43.8 \\
Máximo & 195 & 82 \\
\end{longtable}
}

\textbf{Observación clave:} La distribución de clicks por sesión está
fuertemente sesgada a la derecha (media 9.8, mediana 6), indicando que
la mayoría de los usuarios realizan navegaciones cortas, mientras que
una minoría explora extensivamente.




\end{document}
